A distinctive feature of \Gweek{} is that its optimiser is derived directly from
the equational theory. The equations of \cref{figure:equations}, validated by
the denotational semantics of \parencite{jones_2026}, serve as rewrite rules:
each rule replaces a term with an equal one, and the soundness of the
transformation follows from the full abstraction theorem of the semantics rather
than from ad hoc argument.

The optimiser operates as a \emph{peephole pass} over the CBPV term produced by
the translation of \cref{subsection:translation}. It traverses the term
bottom-up, applying rewrites at each node until a fixed point is reached.
Because the machine uses de Bruijn indices, the rewrites require shifting and
substitution on indices, which we implement as straightforward recursive
traversals. \Cref{figure:rewrites} lists all rewrite rules alongside the
equations that justify them. We now describe them in detail.

\begin{figure}
\begin{small}
\begin{tabular}{@{}lll@{}}
  \toprule
  \textbf{Rewrite} & \textbf{Equation} & \textbf{Law} \\
  \midrule
  \textsc{bind-eta}
    & $\CTerm{\CTo{\Return{V}}{x}{M}} \leadsto \CTerm{\CSubst{M}{V}{x}}$
    & $\beta$-bind \\
  \textsc{dead-bind}
    & $\CTerm{\CTo{\Return{V}}{x}{N}} \leadsto \CTerm{N}$ \quad ($x \notin \textsf{fv}(N)$)
    & $\beta$-bind \\
  \textsc{bind-assoc}
    & $\CTerm{\CTo*{\CTo{M}{x}{N}}{y}{P}} \leadsto \CTerm{\CTo{M}{x}{(\CTo{N}{y}{P})}}$
    & seq.\ assoc. \\
  \textsc{app-bind}
    & $\CTerm{\CApp*{\CTo{M}{x}{N}}{V}} \leadsto \CTerm{\CTo{M}{x}{\CApp{N}{V}}}$
    & seq.\ app. \\
  \textsc{lam-beta}
    & $\CTerm{\CApp*{\Lam{x}{M}}{V}} \leadsto \CTerm{\CSubst{M}{V}{x}}$
    & $\beta$-app \\
  \textsc{force-beta}
    & $\CTerm{\Force*{\Thunk{M}}} \leadsto \CTerm{M}$
    & $\beta$-thunk \\
  \textsc{ifz-beta}
    & $\CTerm{\CIfz{\Zero}{M}{n}{N}} \leadsto \CTerm{M}$ (etc.)
    & $\beta$-ifz \\
  \textsc{match-beta}
    & $\CTerm{\CMatch{\Cons{V}{W}}{M}{x}{xs}{N}} \leadsto \CTerm{\SubstTwo{N}{V}{W}{x}{xs}}$ (etc.)
    & $\beta$-match \\
  \textsc{case-beta}
    & $\CTerm{\CCase{\In[1]{V}}{x}{M}{y}{N}} \leadsto \CTerm{\CSubst{M}{V}{x}}$ (etc.)
    & $\beta$-case \\
  \textsc{dead-end}
    & $\CTerm{\CTo{M}{x}{\Fail}} \leadsto \CTerm{\Fail}$ \quad ($M$ total)
    & dead-end \\
  \textsc{lam-choice}
    & $\CTerm{\Lam*{x}{\Choice{M}{N}}} \leadsto \CTerm{\Choice{\Lam{x}{M}}{\Lam{x}{N}}}$
    & $\lambda$-$\talloblong$ \\
  \textsc{lam-exists}
    & $\CTerm{\Lam*{x}{\Exists{z}[\sigma]{M}}} \leadsto \CTerm{\Exists*{z}[\sigma]{\Lam{x}{M}}}$
    & $\lambda$-$\exists$ \\
  \textsc{lam-equate}
    & $\CTerm{\Lam*{x}{\Equate{V}{W}{M}}} \leadsto \CTerm{\Equate{V}{W}{\Lam{x}{M}}}$
    & $\lambda$-eq \\
  \textsc{eq-choice}
    & $\CTerm{\Equate{V}{W}{\Choice{M}{N}}} \leadsto \CTerm{\Choice{(\Equate{V}{W}{M})}{(\Equate{V}{W}{N})}}$
    & eq-$\talloblong$ \\
  \textsc{eq-exists}
    & $\CTerm{\Equate{V}{W}{\Exists{z}[\sigma]{M}}} \leadsto \CTerm{\Exists*{z}[\sigma]{\Equate{V}{W}{M}}}$
    & eq-$\exists$ \\
  \textsc{cycle}
    & $\CTerm{\Equate{V}{\Succ{V}}{M}} \leadsto \CTerm{\Fail}$ (etc.)
    & occurs \\
  \bottomrule
\end{tabular}
\end{small}
\caption{Peephole rewrite rules and their justifying equations. All equations
  appear in \cref{figure:equations} or are validated by the denotational
  semantics of \parencite{jones_2026}. In addition, the choice identity laws
  ($\CTerm{\Choice{\Fail}{M}} = \CTerm{M}$), the idempotence law
  ($\CTerm{\Choice{M}{M}} = \CTerm{M}$), the reflexivity of unification
  ($\CTerm{\Equate{V}{V}{M}} = \CTerm{M}$), and the narrowing
  decomposition rules are applied.}
\label{figure:rewrites}
\end{figure}

\paragraph{$\beta$-reductions}
The $\beta$-laws of \cref{figure:equations} are the workhorses of the
optimiser. The bind-return law
$\CTerm{\CTo{\Return{V}}{x}{M}} = \CTerm{\CSubst{M}{V}{x}}$
(\textsc{bind-eta}) is the most frequently applied: it eliminates the overhead
of a continuation push and pop by substituting the value directly. A special
case arises when the returned value is a variable
($\CTerm{\CTo{\Return{x}}{y}{M}}$), which is reduced by aliasing $y$ to $x$.
When $x$ does not appear free in the continuation, the bind is dead and the
continuation is returned unchanged (\textsc{dead-bind}).

The remaining $\beta$-laws reduce eliminators applied to known constructors:
\textsc{ifz-beta} reduces $\CTerm{\CIfz{\Zero}{M}{n}{N}}$ to $\CTerm{M}$ and
$\CTerm{\CIfz{\Succ{V}}{M}{n}{N}}$ to $\CTerm{\CSubst{N}{V}{n}}$;
\textsc{match-beta} reduces
$\CTerm{\CMatch{\Cons{V}{W}}{M}{x}{xs}{N}}$ to
$\CTerm{\SubstTwo{N}{V}{W}{x}{xs}}$;
\textsc{case-beta} reduces
$\CTerm{\CCase{\In[1]{V}}{x}{M}{y}{N}}$ to $\CTerm{\CSubst{M}{V}{x}}$.
The \textsc{force-beta} rule reduces $\CTerm{\Force*{\Thunk{M}}}$ to
$\CTerm{M}$, and \textsc{lam-beta} reduces
$\CTerm{\CApp*{\Lam{x}{M}}{V}}$ to $\CTerm{\CSubst{M}{V}{x}}$.

\paragraph{Sequencing laws}
The associativity law (\textsc{bind-assoc}) reassociates
$\CTerm{\CTo*{\CTo{M}{x}{N}}{y}{P}}$ to
$\CTerm{\CTo{M}{x}{(\CTo{N}{y}{P})}}$, exposing opportunities for
\textsc{bind-eta}. The application-bind law (\textsc{app-bind}) rewrites
$\CTerm{\CApp*{\CTo{M}{x}{N}}{V}}$ to $\CTerm{\CTo{M}{x}{\CApp{N}{V}}}$.

\paragraph{$\lambda$-laws}
The optimiser distributes $\lambda$-abstractions over effects using the three
$\lambda$-laws of \cref{figure:equations}: \textsc{lam-choice} rewrites
$\CTerm{\Lam*{x}{\Choice{M}{N}}}$ to
$\CTerm{\Choice{\Lam{x}{M}}{\Lam{x}{N}}}$; \textsc{lam-exists} rewrites
$\CTerm{\Lam*{x}{\Exists{z}[\sigma]{M}}}$ to
$\CTerm{\Exists*{z}[\sigma]{\Lam{x}{M}}}$; and \textsc{lam-equate} rewrites
$\CTerm{\Lam*{x}{\Equate{V}{W}{M}}}$ to
$\CTerm{\Equate{V}{W}{\Lam{x}{M}}}$ (when $\VTerm{V}$ and $\VTerm{W}$ do not
reference $x$). These normalise terms so that effects appear at the outermost
level.

\paragraph{Effect equations}
The identity law $\CTerm{\Choice{\Fail}{M}} = \CTerm{M}$ removes dead branches,
and the idempotence law $\CTerm{\Choice{M}{M}} = \CTerm{M}$ collapses
duplicates. Nested choices are flattened using associativity. The dead-end law
(\textsc{dead-end}) rewrites $\CTerm{\CTo{M}{x}{\Fail}}$ to $\CTerm{\Fail}$
when $\CTerm{M}$ is total---a condition we check by a conservative syntactic
analysis. Reflexive unification $\CTerm{\Equate{V}{V}{M}} = \CTerm{M}$
eliminates trivially satisfied constraints.

\paragraph{Equational constraint rules}
The \textsc{eq-choice} rule distributes an equational constraint over a choice:
$\CTerm{\Equate{V}{W}{\Choice{M}{N}}} =
\CTerm{\Choice{(\Equate{V}{W}{M})}{(\Equate{V}{W}{N})}}$.
The \textsc{eq-exists} rule commutes an equational constraint past an
existential: $\CTerm{\Equate{V}{W}{\Exists{z}[\sigma]{M}}} =
\CTerm{\Exists*{z}[\sigma]{\Equate{V}{W}{M}}}$.
The narrowing equations decompose unification on known constructors: for
example, $\CTerm{\Equate{\Succ{V}}{\Succ{W}}{M}}$ reduces to
$\CTerm{\Equate{V}{W}{M}}$, and
$\CTerm{\Equate{\Cons{V_1}{W_1}}{\Cons{V_2}{W_2}}{M}}$ reduces to
$\CTerm{\Equate{V_1}{V_2}{\Equate{W_1}{W_2}{M}}}$.
Mismatched constructors ($\CTerm{\Equate{\Succ{V}}{\Zero}{M}}$,
$\CTerm{\Equate{\Cons{V}{W}}{\Nil}{M}}$, etc.)\ reduce to $\CTerm{\Fail}$.
The \textsc{cycle} rule detects cyclic constraints
($\CTerm{\Equate{V}{\Succ{V}}{M}}$, $\CTerm{\Equate{W}{\Cons{V}{W}}{M}}$,
etc.)\ and reduces them to $\CTerm{\Fail}$.

\paragraph{Discussion}
The key point is not the magnitude of the improvement but its
\emph{trustworthiness}: because every rewrite is a validated equation, the
optimiser cannot change the observable behaviour of the program. This is a
stronger guarantee than is typical for compiler optimisations, and it is a
direct consequence of basing the implementation on the algebraic theory.
